\documentclass[12pt]{article}
\usepackage[utf8]{inputenc}
\usepackage[T1]{fontenc}
\usepackage[margin=1in]{geometry}
\usepackage{amsmath,amssymb}
\usepackage{array}
\usepackage{enumitem}
\usepackage{fancyhdr}
\usepackage{graphicx}
\usepackage{booktabs}
\usepackage{lastpage}

\pagestyle{fancy}
\fancyhf{}
\fancyhead[L]{\small VIETNAM NATIONAL UNIVERSITY -- HCMC \\ \small UNIVERSITY OF SCIENCE}
\fancyhead[R]{\small \textbf{FINAL EXAMINATION} \\ \small Semester II -- Academic Year 2025--2026}
\fancyfoot[L]{\small (This test has \pageref{LastPage} pages)}
\fancyfoot[R]{\small Page \thepage/\pageref{LastPage}}
\renewcommand{\headrulewidth}{0.4pt}

\setlength{\headheight}{30pt}
\setlength{\parindent}{0pt}
\setlength{\parskip}{6pt}

\begin{document}

\begin{center}
\textbf{\large TEST 3}\\[2pt]
\textbf{Course:} Applied Statistics for Engineering and Scientists 1 \quad
\textbf{Course ID:} STAT451\\[2pt]
\textbf{Duration:} 90 minutes \hspace{1.5cm} \textbf{Materials:} One A4 formula sheet allowed
\end{center}

\vspace{-6pt}
\noindent\textbf{Student name:} \dotfill \quad \textbf{Student ID:} \dotfill \quad \textbf{No.:} \dotfill

\vspace{4pt}
\noindent\textbf{Note:}
\begin{itemize}[leftmargin=*,topsep=2pt,itemsep=0pt]
  \item All calculations must be rounded to \textbf{four} decimal places.
  \item Standard normal $\mathcal{N}(0,1)$ and Student $t$ distribution tables are provided.
\end{itemize}

%--------------------------------------------------------------------
\textbf{Question 1.} (2.0 \emph{points})
In the production of a certain pharmaceutical tablet, the active
ingredient content (in mg) is approximately normally distributed with
mean $\mu = 250$ mg and standard deviation $\sigma = 6$ mg. Tablets
are considered acceptable only if the active-ingredient content is
between 240 mg and 260 mg.
\begin{enumerate}[label=(\alph*),leftmargin=*]
  \item Compute the probability that a randomly chosen tablet is
        \emph{rejected} (i.e., outside specification).
  \item A large batch of 200 tablets is randomly selected from
        production. Using the normal approximation to the binomial
        distribution, calculate the probability that the batch
        contains between 15 and 25 rejected tablets (inclusive).
  \item A quality inspector randomly selects $n = 40$ tablets.
        Determine the probability that the sample mean
        active-ingredient content exceeds 252 mg.
\end{enumerate}

%--------------------------------------------------------------------
\textbf{Question 2.} (2.0 \emph{points})
The impact strength (in foot-pounds) of insulating materials used in
printed-circuit boards is being monitored. Historical records indicate
that impact strength is approximately normally distributed with
population standard deviation $\sigma = 0.45$ foot-pounds. A random
sample of $n = 36$ specimens produced a sample mean of
$\bar{x} = 1.58$ foot-pounds.
\begin{enumerate}[label=(\alph*),leftmargin=*]
  \item Construct a 99\% confidence interval for the true mean
        impact strength $\mu$.
  \item In a separate inspection, 42 out of 180 randomly selected
        circuit boards failed a thermal-shock test. Construct a
        90\% confidence interval for the true proportion $p$ of
        boards that fail the test.
  \item The quality engineer desires a 95\% confidence interval for
        $\mu$ with a maximum error of 0.10 foot-pounds. What
        minimum sample size is required?
\end{enumerate}

%--------------------------------------------------------------------
\textbf{Question 3.} (2.0 \emph{points})
A brochure claims that the mean tensile strength of a certain synthetic
fiber is 84 cN/tex. To test this, a quality control laboratory
measures a random sample of $n = 12$ fibers. Assuming the tensile
strength is normally distributed, the sample yielded
$\bar{x} = 82.1$ cN/tex and $s = 3.4$ cN/tex.
\begin{enumerate}[label=(\alph*),leftmargin=*]
  \item At the significance level $\alpha = 0.05$, test
        $H_0: \mu = 84$ against $H_1: \mu < 84$.
        Clearly state the rejection region, the value of the test
        statistic, and the conclusion in plain English.
  \item A field study involving 500 fibers found that 85 of them
        broke under the specified working load. The manufacturer
        claims that the true breakage rate is less than 20\%.
        Test this claim at $\alpha = 0.05$.
  \item Compute the $p$-value for part (b) and interpret its meaning
        in context.
\end{enumerate}

%--------------------------------------------------------------------
\textbf{Question 4.} (1.5 \emph{points})
A dietician wishes to evaluate a new weight-loss supplement. Ten
randomly selected volunteers were weighed (in kilograms) before
starting the program and again after 8 weeks.
\begin{center}
\begin{tabular}{l|cccccccccc}
Subject & 1 & 2 & 3 & 4 & 5 & 6 & 7 & 8 & 9 & 10 \\ \hline
Before  & 86.4 & 79.2 & 92.1 & 74.8 & 88.5 & 81.0 & 95.3 & 77.6 & 83.9 & 90.7 \\
After   & 83.1 & 77.8 & 88.4 & 73.5 & 85.2 & 79.6 & 91.7 & 76.9 & 81.0 & 86.8 \\
\end{tabular}
\end{center}
Assume the differences are approximately normal.
\begin{enumerate}[label=(\alph*),leftmargin=*]
  \item Test at $\alpha = 0.05$ whether the supplement produces a
        mean weight loss of more than 2 kg. Use the paired $t$-test
        and state all steps.
  \item Construct a 95\% confidence interval for the mean weight
        loss $\mu_D$.
\end{enumerate}

%--------------------------------------------------------------------
\textbf{Question 5.} (2.5 \emph{points})
The final exam score ($y$) and the number of hours spent studying per
week ($x$) are recorded for $n = 9$ randomly chosen students in an
introductory statistics course:
\begin{center}
\begin{tabular}{l|ccccccccc}
$x$ & 2 & 3 & 4 & 6 & 7 & 8 & 10 & 12 & 14 \\ \hline
$y$ & 55 & 58 & 63 & 70 & 72 & 78 & 83 & 88 & 92 \\
\end{tabular}
\end{center}
\begin{enumerate}[label=(\alph*),leftmargin=*]
  \item Determine the least-squares regression line
        $\hat{y} = \hat{\beta}_0 + \hat{\beta}_1 x$. Interpret the
        estimated slope and the intercept in context.
  \item Compute the sample correlation coefficient $r$ between study
        hours and exam score, and interpret its value.
  \item Compute the coefficient of determination $R^{2}$ and
        interpret it.
  \item Using the regression line, predict the exam score for a
        student who studies 9 hours per week and for a student who
        studies 25 hours per week. Comment on the reliability of
        each prediction.
\end{enumerate}

\vspace{10pt}
\begin{center}
\textbf{--- END OF TEST 3 ---}
\end{center}

\end{document}
